\documentclass[11pt,a4paper]{article}

\usepackage[margin=1.05in]{geometry}
\usepackage{amsmath,amssymb,amsthm}
\usepackage{mathtools}
\usepackage{stmaryrd}
\usepackage{microtype}
\setlength{\emergencystretch}{2.5em}
\usepackage{enumitem}
\usepackage{booktabs}
\usepackage{array}
\usepackage{tikz-cd}
\usepackage{xcolor}
\usepackage[colorlinks=true,linkcolor=blue!45!black,urlcolor=blue!60!black,citecolor=blue!45!black]{hyperref}

\newtheorem{theorem}{Theorem}[section]
\newtheorem{proposition}[theorem]{Proposition}
\newtheorem{lemma}[theorem]{Lemma}
\newtheorem{corollary}[theorem]{Corollary}
\newtheorem{conjecture}[theorem]{Conjecture}
\theoremstyle{definition}
\newtheorem{definition}[theorem]{Definition}
\newtheorem{example}[theorem]{Example}
\theoremstyle{remark}
\newtheorem{remark}[theorem]{Remark}
\newtheorem{warning}[theorem]{Warning}
\newtheorem{question}[theorem]{Question}
\newtheorem{thesisclaim}[theorem]{Thesis}

\newcommand{\FF}{\mathbb{F}}
\newcommand{\NN}{\mathbb{N}}
\newcommand{\ZZ}{\mathbb{Z}}
\newcommand{\QQ}{\mathbb{Q}}
\newcommand{\RR}{\mathbb{R}}
\newcommand{\CC}{\mathbb{C}}
\newcommand{\cat}[1]{\mathbf{#1}}
\newcommand{\Set}{\cat{Set}}
\newcommand{\Cont}{\cat{Cont}}
\newcommand{\Poly}{\cat{Poly}}
\newcommand{\Lam}{\Lambda}
\newcommand{\AW}{\mathrm{AW}(3)}
\newcommand{\Weyl}{\mathcal{A}_1}
\newcommand{\Weylq}{\mathcal{A}_1^{q}}
\newcommand{\con}{\mathrel{\vartriangleright}}
\newcommand{\norm}[1]{\lvert #1 \rvert}
\newcommand{\dd}{\,d}
\newcommand{\Had}{\odot}
\DeclareMathOperator{\GL}{GL}
\DeclareMathOperator{\Aut}{Aut}
\DeclareMathOperator{\List}{List}
\DeclareMathOperator{\Bag}{Bag}
\newcommand{\qbin}[2]{\genfrac{[}{]}{0pt}{}{#1}{#2}_{q}}

\definecolor{okgreen}{rgb}{0.0,0.45,0.15}
\definecolor{nored}{rgb}{0.65,0.0,0.0}
\definecolor{amber}{rgb}{0.70,0.45,0.0}
\newcommand{\OK}{\textcolor{okgreen}{\textbf{verified}}}
\newcommand{\NO}{\textcolor{nored}{\textbf{fails}}}
\newcommand{\MAYBE}{\textcolor{amber}{\textbf{open}}}

\title{\bfseries The Organizing Theme Andrews--Askey--Roy Does Not Name\\[0.4em]
\large A research programme: the Weyl algebra as the spine of the special-function
pile, and containers as its categorification}
\author{}
\date{}

\begin{document}
\maketitle

\begin{abstract}
A standard special-functions reference reads as a jumbled pile because it presents the
union of four mutually incompatible organizing themes and commits to none: the Askey
scheme (order by parameter degeneration), monodromy (order by singularities on
$\mathbb{P}^1$), Lie theory (order by which group's representations), and the
holonomic/$D$-module view (order by rank). This note argues for the fourth, on the
ground that it is the only one of which we already possess a categorification. Its
central relation, $[\partial, x] = 1$, is exactly $DU - UD = I$ on Young's lattice ---
Stanley's differential-poset axiom, the corollary targeted by
\texttt{lascoux/lean-plan.md} step 1.8, and the identity that Fomin's local rules
prove bijectively. The $q$-layer closes the loop with the rest of
\texttt{q-bool/}: the $q$-Weyl relation $D_q M_x - q\,M_x D_q = 1$ is the Jackson row
already present in the decategorification table of \texttt{q-container.tex}. The
programme is therefore to climb the Askey scheme from its bottom --- Hermite, the Fock
representation, $\Weyl$ --- towards its top --- Askey--Wilson, Zhedanov's $\AW$ ---
carrying a container model up with us, and to stop honestly wherever the carrying
fails. Five phases, each with a deliverable and a falsification condition. Phase C is
flagged in advance as the likely point of failure, and \S6 says why.
\end{abstract}

\tableofcontents

\section{The complaint, and the diagnosis}

The complaint is that a comprehensive special-functions reference --- Andrews, Askey
and Roy \cite{AAR} is the one at hand --- presents a pile. Chapter after chapter of
families, each with its own orthogonality, recurrence, integral representation and
generating function, and no evident reason why these families and not others.

The diagnosis is not that the book is disorganised. It is that \emph{four} organizing
themes exist, they do not agree, and a comprehensive book presents their union.

\begin{center}
\small
\renewcommand{\arraystretch}{1.3}
\begin{tabular}{@{}llll@{}}
\toprule
theme & orders the pile by & the chart & canonical source \\
\midrule
Askey scheme & degeneration of parameters & a poset of limits & \cite{KLS} \\
monodromy & singularities on $\mathbb{P}^1$ & confluence of punctures & \cite{WW} \\
Lie theory & which group's irreps & group contraction & \cite{Miller,Vilenkin} \\
holonomic / $D$-module & rank over the Weyl algebra & \emph{a finiteness theorem} & \cite{Zeil90,SST} \\
\bottomrule
\end{tabular}
\end{center}

\noindent
They genuinely disagree. The Askey scheme orders by what happens as a parameter tends
to a limit; monodromy orders by how many punctures the differential equation has;
Lie theory orders by the dimension and type of a symmetry group; holonomy orders by a
rank, and produces not a chart but a decision procedure. Any two of these cut the pile
differently, which is why no two books' tables of contents agree.

\begin{remark}[the cheap fix, stated first]\label{rem:cheap}
If the immediate want is relief rather than a programme: the Askey scheme is a
literal one-page directed graph, and Koekoek--Lesky--Swarttouw \cite{KLS} is organized
\emph{as} the chart rather than around it, with the $q$-scheme as a parallel chart
that \cite{AAR} does not have. Thirteen classical families become one poset with
Askey--Wilson on top and Hermite at the bottom. Half an hour, and the pile is a poset.
That is worth doing before anything below.
\end{remark}

\section{The thesis}

\begin{thesisclaim}[\emph{ours}]\label{thesis}
Of the four themes, the holonomic one is the organizing theme to adopt, because it is
the only one whose central object we already possess a categorification of. Its
central object is the Weyl algebra; our categorification of the Weyl algebra is
Fomin's local rule, read as a vertical container isomorphism.
\end{thesisclaim}

The remainder of \S2 establishes the three links this depends on. None of the three is
new; what is new, as far as we can tell, is putting them in one line.

\subsection{Link 1: $DU - UD = I$ is the Weyl algebra}

\begin{definition}\label{def:weyl}
The (first) Weyl algebra is $\Weyl = \CC\langle x, \partial\rangle$ subject to
$\partial x - x\partial = 1$.
\end{definition}

\begin{proposition}[standard]\label{prop:young-weyl}
Let $\Lam = \CC[p_1,p_2,\dots]$ be the ring of symmetric functions and identify
$\CC[\text{Young}] \cong \Lam$ by $\lambda \mapsto s_\lambda$. Then
\[
  U = \text{multiplication by } p_1, \qquad D = \partial/\partial p_1
\]
realise the up and down operators of Young's lattice --- $U s_\lambda = \sum_{\mu \gtrdot
\lambda} s_\mu$ by the Pieri rule, and $D$ is its adjoint for the Hall inner product
--- and satisfy $DU - UD = I$. This is Stanley's differential-poset axiom \cite{Stanley88}.
\end{proposition}

\begin{remark}[a multiplicity space, and a question]\label{rem:mult}
Be precise about what representation this is. As an $\Weyl$-module,
$\Lam = \CC[p_1]\otimes\CC[p_2,p_3,\dots]$, with $\Weyl$ acting on the first factor
only. So Young's lattice carries the Fock representation \emph{with multiplicity
space} $\CC[p_2,p_3,\dots]$; the higher power sums are spectators. The polynomial ring
$\CC[x]$ with $x, d/dx$ --- whose orthogonal basis for the Gaussian measure is the
Hermite polynomials, the bottom of the Askey scheme --- is the same representation with
trivial multiplicity. Hence Question~\ref{q:mult} below.
\end{remark}

\subsection{Link 2: Fomin's rule is the bijective proof of that relation}

Counting the fibres of the vertical container isomorphism of
\texttt{hypergeometric.tex} \S5.2 gives
\[
  \norm{U\alpha \cap U\beta} = \delta_{\alpha\beta} + \norm{D\alpha \cap D\beta},
\]
which is $DU - UD = I$ coefficientwise. The $\delta_{\alpha\beta}$ is the growth-diagram
label $x = 1$ being legal at exactly one face shape. This is step 1.8 of
\texttt{lascoux/lean-plan.md}, and steps 1.3--1.5 are the rule, its reverse, and the
two round trips; 1.3 is complete.

\begin{remark}\label{rem:bijective}
The content of Link 2 is that the relation of Definition~\ref{def:weyl} has a
\emph{bijective} proof, uniform in $(\alpha,\beta)$, in which the two sides are fibres
of one container and the relation is an isomorphism between them rather than an
equation between numbers. That is what makes Thesis~\ref{thesis} more than a
restatement: we do not merely know the relation, we have a morphism that witnesses it.
\end{remark}

\subsection{Link 3: holonomic means finite rank over the Weyl algebra}

\begin{proposition}[reported, \cite{Zeil90}]\label{prop:holonomic}
A function is \emph{holonomic} when it is annihilated by a left ideal of the Weyl
algebra $\mathcal{A}_n$ of maximal possible dimension --- equivalently, when the
$\mathcal{A}_n$-module it generates has finite holonomic rank. Holonomic functions are
closed under sum, product, and definite summation or integration over a variable;
identities between them are decidable, and the Wilf--Zeilberger machinery is the
decision procedure.
\end{proposition}

So this theme does not produce a chart. It produces a finiteness statement: the pile is
not a pile but the set of finite-rank modules over one algebra, and the apparent
profusion of families is the profusion of ideals, not of phenomena. The reason for
preferring it here is Thesis~\ref{thesis}: \emph{that algebra} is the one we can
categorify.

\subsection{The square}

\[
\begin{tikzcd}[row sep=2.0em, column sep=3.2em]
\Cont \text{ (vertical isos)} \arrow[r, "\text{count fibres}"] \arrow[d, "q\text{-deform}"']
  & \Weyl\text{-modules} \arrow[r, "\text{generating function}"] \arrow[d, "q\text{-deform}"]
  & \text{special functions} \arrow[d, "q\text{-deform}"] \\
\Cont_q \arrow[r, "\text{count fibres}"]
  & \Weylq\text{-modules} \arrow[r]
  & q\text{-special functions}
\end{tikzcd}
\]

\noindent
The top-left corner is \texttt{lascoux/}; the top-middle is Stanley and Zeilberger; the
top-right is \cite{AAR}. The bottom row is \texttt{q-bool/}. Every arrow in the top row
exists. The claim of the programme is that the rest of the square can be filled in, and
the phases of \S4 are the attempts.

\section{The $q$-layer is already half built}

\begin{proposition}\label{prop:qweyl}
Let $D_q f(x) = \dfrac{f(qx) - f(x)}{(q-1)x}$ be the Jackson derivative and $M_x$
multiplication by $x$. Then
\[
  D_q M_x - q\, M_x D_q = 1 .
\]
\end{proposition}

\begin{proof}
$D_q(xf) = \dfrac{qx f(qx) - x f(x)}{(q-1)x} = \dfrac{q f(qx) - f(x)}{q-1}$ and
$x D_q f = \dfrac{f(qx)-f(x)}{q-1}$, so
$D_q(xf) - q\,x D_q f = \dfrac{q f(qx) - f(x) - q f(qx) + q f(x)}{q-1} = f(x)$. \OK
\end{proof}

\begin{corollary}[\emph{ours}]\label{cor:jackson-row}
The decategorification table of \texttt{q-container.tex} contains the row
``derivative $\partial$ (hyperplanes) $\leftrightarrow$ Jackson $D_q$'' and the row
``$\mathrm{Ind}_*,\mathrm{Res}_*$ $\leftrightarrow$ multiplication, comultiplication''.
Together these are the two generators of Proposition~\ref{prop:qweyl}. So a container
model of the $q$-Weyl algebra is not a new construction to be invented: its generators
are already in the table, and what is missing is the \emph{relation} --- the
$q$-analogue of Link 2.
\end{corollary}

This is the single most encouraging fact in this note, and it is what makes Phase D
cheap rather than speculative.

\subsection{The top of the scheme}

\begin{remark}[reported]\label{rem:AW}
The Askey scheme has a top as well as a bottom. Askey--Wilson polynomials sit above
every other family in \cite{KLS}, and the algebra governing them is Zhedanov's
Askey--Wilson algebra $\AW$ \cite{Zhedanov,Terwilliger} --- three generators with
$q$-commutator relations, degenerating down the scheme as the parameters degenerate.
\end{remark}

So the programme has a shape: $\Weyl$ at the bottom (Hermite, where we have a
categorification), $\AW$ at the top (Askey--Wilson, where we have none), and the Askey
scheme as the path between. Climb it carrying the container model, and record where the
model stops being carried.

\section{The programme}

Five phases. Each states its deliverable, its dependencies, an honest confidence that
it will work, and --- the part that matters --- what observation would kill it.

\subsection*{Phase 0 --- orientation (days)}

\begin{enumerate}[nosep,label=\textbf{0.\arabic*}]
  \item Read the chart in \cite{KLS}, per Remark~\ref{rem:cheap}. Deliverable: the
        poset, drawn, with the $q$-scheme beside it.
  \item Read \cite{Zeil90}. Twenty pages, and it is the thesis statement for the
        chosen theme.
  \item Settle Question~3.9 of \texttt{hypergeometric.tex}: is the class of
        $q$-containers closed under the Hadamard product, pulled back along $\dim$?
        Cheap, independent of everything else, and needed by Phase D.
\end{enumerate}

\noindent\textbf{Confidence:} n/a, this is reading. \textbf{Cost:} under a week.

\subsection*{Phase A --- the Weyl algebra at the container level}

\textbf{Deliverable.} Finish \texttt{lascoux/} Phase 1 (steps 1.4--1.8): the reverse
local rule, both round trips, the assembly as a vertical morphism in $\Cont$, and
$DU - UD = I$ as a corollary with no \texttt{sorry}. Then the harder half: a container
model of $\Weyl$ \emph{as an algebra}, not merely of its defining relation ---
objects for $x$ and $\partial$, a composition, and the relation as a specified
isomorphism.

\textbf{Dependencies.} None. Step 1.3 is complete.

\textbf{Confidence.} Steps 1.4--1.8: \textbf{high} (85\%). The reverse rule is the
forward rule's six cases run backwards, and the only case needing new work is the
bounce, where \texttt{diffRow} needs a counterpart that finds the row to \emph{remove}.
The algebra model: \textbf{medium} (50\%), because composition of vertical morphisms
over different shape pairs is exactly the composability problem that thesis \S4.2
records as the optic obstruction.

\textbf{Falsified if} the round trips of step 1.5 fail --- i.e.\ if the rule as stated
is not a bijection on the nose but only up to something. That would not be a disaster
(it would make the structure a lens rather than an iso, which \S4.2 half expects) but
it would change every statement downstream from ``isomorphism'' to ``optic''.

\subsection*{Phase B --- WZ certificates as container morphisms}

\textbf{Deliverable.} An answer to Question~5.10 of \texttt{hypergeometric.tex}: is a
Wilf--Zeilberger certificate a vertical morphism in $\Cont$ over the face data of a
two-dimensional growth diagram? Calibration case: $q$-Vandermonde, laid out as a
growth diagram rather than as a recurrence, whose bijection is already known
(the subspace-intersection argument, \texttt{hypergeometric.tex} Cor.~6.6). The test is
sharp: \emph{does the WZ certificate for $q$-Vandermonde reproduce the subspace
bijection?}

\textbf{Dependencies.} Phase A for the notion of vertical morphism; Phase 0.2.

\textbf{Confidence.} \textbf{Medium-low} (35\%) in general, \textbf{high} (75\%) for
the calibration case alone. The general obstruction is that WZ certificates are
rational functions and need not take integer values, so there need be no bijection to
find. The reason to try anyway is the payoff: if it works, a decision procedure becomes
a machine for producing container isomorphisms.

\textbf{Falsified if} the $q$-Vandermonde certificate is integral but demonstrably not
the subspace bijection, \emph{and} no reindexing relates them. Then WZ and Fomin are
cousins and \S5 of \texttt{hypergeometric.tex} should end at its Warning~5.7.

\subsection*{Phase C --- the Askey scheme as a poset of containers}

\textbf{Deliverable.} A container-level meaning for the \emph{edges} of the Askey
scheme. Each edge is a parameter limit; the question is what a limit of containers is.

\textbf{Dependencies.} Phases A and B; Phase 0.1.

\textbf{Confidence.} \textbf{Low} (15\%), and this is flagged in advance as the likely
point of failure of the whole programme. See \S6.

\textbf{Falsified if} --- as expected --- the limits are genuinely analytic, i.e.\ the
degenerations exist only after passing to series with real parameters. The honest
outcome is then a sub-poset: the part of the Askey scheme whose edges \emph{are}
container constructions, with the rest marked as outside.

\subsection*{Phase D --- the $q$-layer}

\textbf{Deliverable.} The $q$-analogue of Link 2: a container-level isomorphism whose
fibre count is Proposition~\ref{prop:qweyl}. By Corollary~\ref{cor:jackson-row} the
generators are already present, so the work is the relation.

\textbf{Dependencies.} Phase A (for the pattern to imitate); Phase 0.3; and
\emph{critically} Question~6.1 of \texttt{list-and-bag.tex} --- does $\Cont_q$ have a
usable $K_0$? --- wherever a minus sign is needed.

\textbf{Confidence.} \textbf{Medium-high} (60\%). Higher than Phase B because the
target is a single relation rather than a general procedure, and because
\texttt{hypergeometric.tex} \S6 argues the $q$-side is the structurally honest one:
there the obligatory factor $(q;q)_n$ is $\norm{\GL_n(\FF_q)}$ up to the torus, i.e.\
an automorphism group order by construction rather than by the accident that
$(1)_n = \norm{S_n}$.

\textbf{Falsified if} $K_0(\Cont_q)$ does not exist usefully. Then the $q$-layer is
confined to statements with no subtraction in them, which excludes the commutator.

\subsection*{Phase E --- Lean}

\textbf{Deliverable.} Phase A formalised (this is already \texttt{lascoux/lean/}); the
$q$-Weyl relation of Proposition~\ref{prop:qweyl} formalised, which is elementary; and
whichever of B and D survive.

\textbf{Dependencies.} Whatever it formalises.

\textbf{Confidence.} \textbf{High} (80\%) for what is in scope, because
\texttt{lascoux/lean-plan.md} \S2 has already inventoried Mathlib and the missing
pieces are known: no Schur polynomials, no RSK, no Hall inner product.

\textbf{Warning.} Proposition~\ref{prop:young-weyl} as stated needs the Hall inner
product, which Mathlib does not have. Either prove $DU - UD = I$ combinatorially on
Young diagrams --- which is what step 1.8 does, and is the reason that step is phrased
as a cardinality and not as an operator identity --- or accept a substantial detour.
Prefer the former. The operator phrasing is for this note, not for the kernel.

\section{Dependency summary}

\begin{center}
\small
\renewcommand{\arraystretch}{1.25}
\begin{tabular}{@{}llll@{}}
\toprule
phase & needs & confidence & status \\
\midrule
0 \ orientation & --- & n/a & not started \\
A \ Weyl at container level & --- & 85\% / 50\% & step 1.3 done \\
B \ WZ as container morphism & A, 0.2 & 35\% / 75\% & \MAYBE \\
C \ Askey scheme as poset & A, B, 0.1 & 15\% & \MAYBE, see \S6 \\
D \ the $q$-layer & A, 0.3, $K_0$ question & 60\% & generators present \\
E \ Lean & the above & 80\% & \texttt{lascoux/lean/} live \\
\bottomrule
\end{tabular}
\end{center}

\section{What would kill the programme, and the one that probably will}

Recorded separately from the phases, because a programme that cannot say how it fails
is not a programme.

\begin{enumerate}[label=\textbf{F\arabic*},leftmargin=2.6em]
  \item \textbf{Degeneration is analytic, not combinatorial.} The edges of the Askey
        scheme are limits: $q \to 1$, a parameter $\to \infty$, a rescaling. Limits are
        not container operations; they are not operations on $\Poly$ at all, which is a
        rig category with no notion of convergence. \emph{This is the expected failure,
        and it is Phase C.} The likely honest outcome is that the container reading
        covers the \emph{vertices} of the scheme (each family, at integer parameters)
        and none of the \emph{edges}, leaving a discrete set where a poset was wanted.
  \item \textbf{No minus signs.} $\Poly$ and $\Cont_q$ are rig categories. A commutator
        $DU - UD$ needs subtraction, so Link 2 at the container level works only because
        it is phrased as $\norm{U \cap U} = \delta + \norm{D \cap D}$ --- all additions.
        Any step that genuinely needs the difference needs
        $K_0$, which is Question~6.1 of \texttt{list-and-bag.tex} and is open.
  \item \textbf{Non-integral certificates.} Phase B's general case, as above.
  \item \textbf{The multiplicity space is a coincidence.} If $\CC[p_2,p_3,\dots]$ in
        Remark~\ref{rem:mult} has nothing to do with the higher parameters of the Askey
        families, then the bottom of the scheme is where the categorification lives and
        there is no ladder to climb.
\end{enumerate}

\begin{warning}[the shape of the honest outcome]\label{warn:honest}
Taking F1 and F4 together, the most likely true statement at the end of this programme
is narrower than Thesis~\ref{thesis}: \emph{the Weyl algebra is categorified by
containers; the Askey scheme is not}. That would still be worth having. It would say
that the organizing theme \cite{AAR} lacks is the holonomic one, that the holonomic
one has a combinatorial foundation, and that the scheme sitting on top of it is an
analytic structure that the combinatorics does not see. Three sentences, all of them
checkable, is a better outcome than an unfalsifiable programme.
\end{warning}

\section{Out of scope, with reasons}

\begin{description}[leftmargin=1.8em,style=nextline,nosep]
\item[Asymptotics and orthogonality measures.] The analytic heart of \cite{AAR}. Every
technique there is estimation, and nothing in \S2 bears on it. Not a gap --- a
different subject.
\item[Painlev\'e and isomonodromy.] The natural continuation of the monodromy theme,
which this note does not adopt. Deliberately left to one side; if Phase C fails as
expected, monodromy is where to look next, because confluence of punctures is a
genuinely structural operation in a way that parameter limits are not.
\item[Multivariable and GKZ.] \cite{SST} organizes multivariable hypergeometric
functions by toric geometry, and the polytope is plausibly a container-like object.
This is the most interesting excluded item and should be revisited after Phase D, not
before.
\item[Macdonald polynomials.] Already scoped out of
\texttt{lascoux/lean-plan.md}, for the reasons given there. The same reasons apply.
\end{description}

\section{Open questions, collected}

\begin{question}\label{q:mult}
Does the multiplicity space $\CC[p_2,p_3,\dots]$ of Remark~\ref{rem:mult} correspond to
the higher parameters of the Askey families? Concretely: is there an assignment of
families in the scheme to subspaces or quotients of $\CC[p_2,p_3,\dots]$ under which
degeneration is a map of multiplicity spaces? If yes, F4 is answered and Phase C has a
mechanism. If no, the programme stops at the Fock representation.
\end{question}

\begin{question}\label{q:aw}
Is there a container model of $\AW$? Equivalently: does Zhedanov's algebra have a
presentation in which the relations are fibre-count identities with no subtraction, as
$DU - UD = I$ does in the form $\norm{U\cap U} = \delta + \norm{D\cap D}$? This is the
top-of-the-scheme question and it is the right one to ask before attempting Phase C,
because a negative answer here closes Phase C immediately and cheaply.
\end{question}

\begin{question}\label{q:lens}
If Phase A's step 1.5 shows the local rule is an optic rather than an isomorphism ---
which thesis \S4.2 half predicts --- what does the \emph{decategorification} of an
optic look like? An isomorphism decategorifies to an equation. A lens decategorifies
to what?
\end{question}

\begin{thebibliography}{99}

\bibitem{AAR} George E.\ Andrews, Richard Askey and Ranjan Roy, \emph{Special
Functions}, Encyclopedia of Mathematics and its Applications 71, Cambridge University
Press, 1999.

\bibitem{KLS} Roelof Koekoek, Peter A.\ Lesky and Ren\'e F.\ Swarttouw,
\emph{Hypergeometric Orthogonal Polynomials and Their $q$-Analogues}, Springer, 2010.
The Askey scheme and the $q$-Askey scheme, as charts.

\bibitem{Miller} Willard Miller Jr., \emph{Lie Theory and Special Functions},
Academic Press, 1968.

\bibitem{SST} Mutsumi Saito, Bernd Sturmfels and Nobuki Takayama, \emph{Gr\"obner
Deformations of Hypergeometric Differential Equations}, Springer, 2000.

\bibitem{Stanley88} Richard P.\ Stanley, \emph{Differential posets}, J.\ Amer.\ Math.\
Soc.\ \textbf{1} (1988), 919--961.

\bibitem{Terwilliger} Paul Terwilliger, \emph{The universal Askey--Wilson algebra},
SIGMA \textbf{7} (2011), 069.

\bibitem{Vilenkin} N.\ Ja.\ Vilenkin, \emph{Special Functions and the Theory of Group
Representations}, AMS Translations of Mathematical Monographs 22, 1968.

\bibitem{WW} E.\ T.\ Whittaker and G.\ N.\ Watson, \emph{A Course of Modern Analysis},
4th ed., Cambridge University Press, 1927. Ch.\ X for the hypergeometric equation and
its singularities.

\bibitem{Zeil90} Doron Zeilberger, \emph{A holonomic systems approach to special
functions identities}, J.\ Comput.\ Appl.\ Math.\ \textbf{32} (1990), 321--368.

\bibitem{Zhedanov} Alexei S.\ Zhedanov, \emph{``Hidden symmetry'' of Askey--Wilson
polynomials}, Theoret.\ and Math.\ Phys.\ \textbf{89} (1991), 1146--1157.

\end{thebibliography}

\end{document}
